\documentclass[12pt]{article}
\usepackage[margin=1in]{geometry}
\usepackage[T1]{fontenc}
\usepackage[utf8]{inputenc}
\usepackage{lmodern}
\usepackage{hyperref}

\title{Model Understanding and Data Exploration Notes}
\author{KERI Project}
\date{\today}

\begin{document}

\maketitle

\section{How My DAG Model Worked (Plain Language for a Medical Audience)}

I built a simple directed acyclic graph (DAG), which was a map of assumed cause-and-effect relationships.
Each dot represented one variable:
\begin{itemize}
	\item The first variable was \textbf{Obesity}.
	\item The second variable was \textbf{Diabetes}.
	\item The third variable was \textbf{Cancer}.
\end{itemize}

I drew one-way arrows between these variables to represent the direction of influence.
Because the graph was acyclic, arrows did not loop back.

I used these arrows:
\begin{itemize}
	\item \textbf{Obesity $\rightarrow$ Diabetes}
	\item \textbf{Obesity $\rightarrow$ Cancer}
	\item \textbf{Diabetes $\rightarrow$ Cancer}
\end{itemize}

I used this structure to ask a focused question:
``If diabetes changed, how much did cancer risk change, after accounting for the pathway through obesity?''

I also tested whether the result was stable using internal checks.
For example, I repeated the analysis under perturbations (such as placebo-style and subset checks) to see whether the effect estimate remained directionally similar.

So, the DAG served as an explicit assumption map.
It helped me separate correlation from a more defensible causal interpretation,
and it made the model logic transparent for review.

\section{Merged NHANES Columns and How I Explored the Data}

The merged file was:
\begin{center}
\texttt{api/nhanes\_data/nhanes\_merged.csv}
\end{center}

The merge script kept original NHANES variable names and added a prefix so I could track each source table.
All tables were joined using \texttt{SEQN}, which was the participant ID key.

\subsection{What the Prefixes Meant}

I organized the merged columns into four main families:
\begin{itemize}
	\item \textbf{DEMO\_...}: Demographics (age, sex, race/ethnicity, income, survey design variables).
	\item \textbf{BMX\_...}: Body measurements (height, weight, BMI, waist, and related body metrics).
	\item \textbf{DIQ\_...}: Diabetes questionnaire responses.
	\item \textbf{MCQ\_...}: Medical conditions questionnaire responses (including cancer-related questions).
\end{itemize}

Examples from the merged header:
\begin{itemize}
	\item \texttt{DEMO\_RIDAGEYR}: age in years.
	\item \texttt{BMX\_BMXBMI}: body mass index.
	\item \texttt{DIQ\_DIQ010}: diabetes status question.
	\item \texttt{MCQ\_MCQ220}: cancer status question used to derive the model target.
\end{itemize}

\subsection{Model-Ready Columns in the Same File}

The preparation step also created three model-ready columns:
\begin{itemize}
	\item \texttt{Obesity}: derived from \texttt{BMX\_BMXBMI} (BMI $\geq 30$ mapped to 1).
	\item \texttt{Diabetes}: derived from \texttt{DIQ\_DIQ010} (Yes/No recoded to 1/0).
	\item \texttt{Cancer}: derived from \texttt{MCQ\_MCQ220} (Yes/No recoded to 1/0).
\end{itemize}

So, the merged dataset contained both raw NHANES-prefixed variables and these binary fields used directly by the DAG model.

\subsection{How I Explored the Data}

I explored the file in a practical sequence:
\begin{enumerate}
	\item I checked the header to see all column names.
	\item I looked at a few rows to confirm values looked reasonable.
	\item I checked missing values to find sparse columns.
	\item I summarized key variables to understand distributions.
	\item I counted class balance for \texttt{Obesity}, \texttt{Diabetes}, and \texttt{Cancer}.
\end{enumerate}

\subsection{Reproducible Exploration Commands}

\textbf{Python (pandas) quick scan:}

\begin{verbatim}
import pandas as pd

df = pd.read_csv("api/nhanes_data/nhanes_merged.csv")

# 1) Shape and first columns
print(df.shape)
print(df.columns[:30].tolist())

# 2) Preview rows
print(df.head(5))

# 3) Missingness (top 20 most missing columns)
missing = df.isna().mean().sort_values(ascending=False)
print(missing.head(20))

# 4) Summary stats for numeric columns
print(df.describe(include=["number"]).T.head(20))

# 5) Class balance for model columns
for col in ["Obesity", "Diabetes", "Cancer"]:
	if col in df.columns:
		print("\n", col)
		print(df[col].value_counts(dropna=False))
\end{verbatim}

\textbf{Header-only check from terminal:}

\begin{verbatim}
python - <<'PY'
import pandas as pd
df = pd.read_csv("api/nhanes_data/nhanes_merged.csv", nrows=0)
print(len(df.columns))
print(df.columns.tolist()[:40])
PY
\end{verbatim}

These steps helped me quickly identify what each column family represented,
which variables had acceptable completeness,
and how DAG inputs were generated from the merged NHANES dataset.

\end{document}
